\section{The counting equation and its integral coefficients}\label{sec:construction}
This section recalls the counting construction with the normalization needed for higher coefficients. The construction is due to the companion counting paper \citep[Theorems~1.1--1.2, Sections~2--3]{Counting}; the calculations below check the integral information that its first Hasse coefficient does not record.

Write a width-three formula in conjunctive normal form (CNF) as a conjunction of local relations $f$, with variable scopes $I_f$ that are sets of distinct variables. Different clauses may have the same scope. A declared variable absent from the formula is given one unary tautological relation. If the incidence set is empty, add a fresh variable with a unary clause forcing it to zero; this makes the incidence set nonempty without changing the count. Write $\mathcal I=\{(f,i):i\in I_f\}$ for that set. For each incidence $e=(f,i)$ introduce $w_{e,0},w_{e,1}$. Put
\[
 P_f=\sum_{a\models f}\prod_{i\in I_f}w_{(f,i),a_i},\qquad
 Q_i=\prod_{f:i\in I_f}w_{(f,i),0}+
     \prod_{f:i\in I_f}w_{(f,i),1},
\]
\[
 P=\prod_fP_f\prod_iQ_i,\qquad
 Z=\prod_e w_{e,0}w_{e,1},\qquad b=2|\mathcal I|.
\]
The coefficient $[Z]P$ is the number $S$ of satisfying assignments: the variable factor chooses a consistent label, and every neighboring clause factor must choose its complement. This gives a bijection, including the multiplicities of equal monomials. In particular $P$ is an integral, homogeneous, positive circuit of degree $b$, and every individual variable has degree at most two.

Introduce $p-1$ tags and set $u=b+p-1$ and
\[
 G=-t_1^{p-1}P+\sum_{j=2}^{p-1}t_j^{p-1}Z.
\]
The empty sum is zero. This is degree $u$ in $u$ variables. Choose the least power $Q=p^r$ with $Q-1>u$ and $Q-1\geq p$, put $m=Q-1$ and $d=Q+1$, and relabel the variables of $G$ as $x_1,\ldots,x_u$. Define integral polynomials
\[
 A=G\prod_{i=u+1}^{m}x_i+x_1^m,\qquad
 F=yzA+y^{Q+1}+\sum_{i=1}^{m}x_i^{Q+1}-x_1z^Q.
\]
Reduction modulo $p$ defines the counting hypersurface $X_\Phi=V(F)\subseteq\PP^{d-1}$.

For clarity, its smoothness is uniform in every coefficient of $G$. In characteristic $p$,
\[
 F_y=zA+y^Q,\quad F_z=yA,\quad yF_y-zF_z=y^{Q+1}.
\]
A common zero of the partial derivatives has $y=0$, then $x_i=0$ for $i>1$ and $x_1^Q=z^Q$. At a projective point both $x_1,z$ are nonzero. The padding makes $A(x_1,0,\ldots,0)=x_1^m$, so $F_y=zx_1^m\ne0$, a contradiction. Adjunction and the hypersurface cohomology sequence give trivial canonical sheaf, a one-dimensional top coherent group, and zero intermediate coherent groups. After every extension of the base field, the same sequence gives $H^0(\mathcal O)$ equal to that field; the smooth geometric hypersurface is therefore connected and integral. These arguments also establish the hypotheses used for the formal group.

Let
\[
 a_N(F)=[(x_1\cdots x_myz)^{N-1}]F^{N-1},\qquad
 T_M(\Phi)=[Z^M]P^M.
\]
\begin{proposition}\label{prop:early-counting-coefficients}
If $p^k-1<Q(Q+1)$ and $M=(p^k-1)/(p-1)$, then over the integers
\[
 a_{p^k}(F)=(-1)^M
 \frac{(p^k-1)!}{(M!)^{p-1}}\,T_M(\Phi).
\]
The displayed multiplier is a $p$-adic unit congruent to one modulo $p$.
\end{proposition}
\begin{proof}
Put $n=p^k-1$. In a contributing term let $a,c,t$ count the selections of $y^{Q+1}$, $-x_1z^Q$, and $yzA$, respectively. The $y,z$ exponents give
$(Q+1)a+t=n=Qc+t$, so $a=Qj,c=(Q+1)j$ for an integer $j\geq0$, and $t=n-Q(Q+1)j$. The assumed inequality forces $j=0$. All $n$ factors therefore select $yzA$. Any padding variable has target exponent $n$, and occurs only in the first summand of $A$, with exponent one. Consequently every factor selects that summand. The coefficient reduces exactly to the diagonal coefficient of $G^n$.
Each tag then forces precisely $M$ selections of its summand. Removing the $p-2$ copies of the $Z$ summand leaves $[Z^M]P^M$, with the asserted sign and multinomial coefficient. The base-$p$ digits of $M$ are $k$ ones. Adding $p-1$ copies produces $k$ digits equal to $p-1$ without a carry, so the multinomial is a unit. Its residue is $((p-1)!)^k=(-1)^k$ by the elementary factorial digit congruence. For odd $p$, $M\equiv k\pmod2$, so the sign cancels this residue. At $p=2$ signs agree modulo $p$.
\end{proof}

